% Gerado automaticamente. Valores medios da fonte experimental; caracteristicos
% obtidos por f_k = 0,70 f_m (compressao e flexao) e f_k = 0,54 f_m (cisalhamento).
\begin{table}[!ht]
\caption{Propriedades mecânicas das 40 espécies de madeira tropical brasileira e enquadramento resultante nas classes de resistência da NBR 7190-3.\label{tab:especies}}
\begin{threeparttable}
\footnotesize
\setlength{\tabcolsep}{3pt}
\begin{tabular*}{\textwidth}{@{\extracolsep{\fill}}rllrrrrrrr@{\extracolsep{\fill}}}
\toprule
\# & Nome popular & Nome científico & $\rho_{ap}$ & $f_{c0,m}$ & $f_{c0,k}$ & $f_{m,m}$ & $f_{v0,m}$ & $E_{M0}$ & Classe \\
 & & & (kg/m\textsuperscript{3}) & (MPa) & (MPa) & (MPa) & (MPa) & (MPa) & (NBR) \\
\midrule
01 & Angelim-amargoso & \textit{Vataireea fusca} & 772 & 60 & 42,0 & 89 & 16,0 & 14947 & D40 \\
02 & Angelim-araroba & \textit{Vataireopsis araroba} & 674 & 50 & 35,0 & 75 & 12,0 & 11457 & D30 \\
03 & Angelim-ferro & \textit{Hymenolobium sp.} & 1163 & 82 & 57,4 & 134 & 19,4 & 19938 & D50 \\
04 & Angelim-pedra-verdadeiro & \textit{Dinizia excelsa} & 1131 & 78 & 54,6 & 110 & 18,9 & 15215 & D50 \\
05 & Angelim-pedra & \textit{Hymenolobium petraeum} & 663 & 58 & 40,6 & 84 & 13,3 & 10755 & D40 \\
06 & Angelim-saia & \textit{Vataireea sp.} & 764 & 63 & 44,1 & 110 & 14,6 & 17561 & D40 \\
07 & Angico-preto & \textit{Piptadenia macrocarpa} & 888 & 73 & 51,1 & 120 & 24,5 & 16498 & D50 \\
08 & Branquilho & \textit{Sebastiania commersoniana} & 810 & 49 & 34,3 & 83 & 16,0 & 15490 & D30 \\
09 & Cafearana & \textit{Andira stipulacea} & 678 & 58 & 40,6 & 95 & 10,2 & 13982 & D40 \\
10 & Canafístula & \textit{Cassia ferruginea} & 860 & 52 & 36,4 & 89 & 18,5 & 14769 & D30 \\
11 & Casca-grossa & \textit{Ocotea odorifera} & 788 & 57 & 39,9 & 107 & 13,4 & 16802 & D30 \\
12 & Castelo & \textit{Calycophyllum multiflorum} & 759 & 55 & 38,5 & 103 & 21,3 & 11375 & D30 \\
13 & Catanudo & \textit{Calophyllum sp.} & 804 & 51 & 35,7 & 83 & 16,4 & 14729 & D30 \\
14 & Cedro-amargo & \textit{Cedrela odorata} & 514 & 38 & 26,6 & 62 & 10,2 & 9176 & D20 \\
15 & Cedro-doce & \textit{Cedrela sp.} & 512 & 33 & 23,1 & 58 & 9,4 & 8866 & D20 \\
16 & Cedrorana & \textit{Cedrelinga catenaeformis} & 566 & 41 & 28,7 & 61 & 11,9 & 10032 & D20 \\
17 & Champanhe & \textit{Dipteryx odorata} & 1090 & 93 & 65,1 & 165 & 17,8 & 25174 & D60 \\
18 & Copaíba & \textit{Copaifera cf. reticulata} & 695 & 50 & 35,0 & 87 & 14,6 & 13572 & D30 \\
19 & Cupiúba & \textit{Goupia glabra} & 839 & 54 & 37,8 & 79 & 17,1 & 13148 & D30 \\
20 & Cutiúba & \textit{Goupia paraensis} & 1152 & 79 & 55,3 & 127 & 17,9 & 16984 & D50 \\
21 & Garapa & \textit{Apuleia leiocarpa} & 920 & 73 & 51,1 & 119 & 19,6 & 16923 & D50 \\
22 & Goiabão & \textit{Planchonella pachycarpa} & 938 & 49 & 34,3 & 107 & 14,1 & 18367 & D30 \\
23 & Guaicará & \textit{Luetzelburgia sp.} & 995 & 66 & 46,2 & 112 & 19,0 & 13866 & D40 \\
24 & Guarucaia & \textit{Peltophorum vogelianum} & 916 & 62 & 43,4 & 96 & 20,3 & 15002 & D40 \\
25 & Ipê & \textit{Tabebuia serratifolia} & 1065 & 78 & 54,6 & 127 & 21,3 & 16670 & D50 \\
26 & Itaúba & \textit{Mezilaurus itauba} & 908 & 69 & 48,3 & 117 & 18,8 & 17345 & D40 \\
27 & Jatobá & \textit{Hymenaea sp.} & 1084 & 91 & 63,7 & 159 & 25,5 & 21367 & D60 \\
28 & Louro-preto & \textit{Ocotea sp.} & 680 & 55 & 38,5 & 85 & 13,8 & 13556 & D30 \\
29 & Maçaranduba & \textit{Manilkara huberi} & 1143 & 83 & 58,1 & 136 & 24,9 & 18184 & D50 \\
30 & Mandioqueira & \textit{Qualea paraensis} & 855 & 71 & 49,7 & 115 & 17,2 & 18679 & D40 \\
31 & Oiticica-amarela & \textit{Clarisia racemosa} & 756 & 70 & 49,0 & 108 & 17,8 & 14491 & D40 \\
32 & Oiuchu & \textit{Pradosia sp.} & 931 & 77 & 53,9 & 123 & 20,1 & 16709 & D50 \\
33 & Parinari & \textit{Parinari excelsa} & 792 & 60 & 42,0 & 112 & 13,4 & 16457 & D40 \\
34 & Piolho & \textit{Tapirira sp.} & 828 & 62 & 43,4 & 76 & 14,7 & 11790 & D40 \\
35 & Quarubarana & \textit{Erisma uncinatum} & 544 & 38 & 26,6 & 67 & 9,6 & 8842 & D20 \\
36 & Rabo-de-arraia & \textit{Vochysia haenkeana} & 729 & 60 & 42,0 & 87 & 13,4 & 14324 & D40 \\
37 & Sucupira & \textit{Diplotropis incexis} & 1103 & 94 & 65,8 & 147 & 20,1 & 20518 & D60 \\
38 & Tachi & \textit{Tachigali myrmecophila} & 1050 & 88 & 61,6 & 140 & 20,9 & 20813 & D60 \\
39 & Tatajuba & \textit{Bagassa guianensis} & 945 & 79 & 55,3 & 110 & 19,9 & 17905 & D50 \\
40 & Umirana & \textit{Qualea retusa} & 705 & 54 & 37,8 & 67 & 14,5 & 10794 & D30 \\
\bottomrule
\end{tabular*}
\begin{tablenotes}\footnotesize
\item $\rho_{ap}$: massa específica aparente a 12\% de umidade; $f_{c0}$: resistência à compressão paralela às fibras; $f_{m}$: resistência à flexão; $f_{v0}$: resistência ao cisalhamento; $E_{M0}$: módulo de elasticidade na flexão estática. Os índices $m$ e $k$ denotam valores médio e característico.
\item Valores característicos obtidos por $f_k = 0{,}70\,f_m$ para compressão e flexão e $f_k = 0{,}54\,f_m$ para cisalhamento, conforme a ABNT NBR 7190-1:2022. \tofill{[CONFERIR contra o artigo-fonte: se os valores caracteristicos ja forem reportados, usar os originais e remover esta nota.]}
\item Fonte: \tofill{[CITAR o artigo do colega com os ensaios]}.
\end{tablenotes}
\end{threeparttable}
\end{table}
